\documentclass[border=8pt]{standalone}
\usepackage{ctex}
\usepackage{amsmath}
\usepackage{tikz}
\begin{document}
\definecolor{mmRootFill}{RGB}{18,85,126}
\definecolor{mmLOneFill}{RGB}{234,148,136}
\begin{tikzpicture}[
  mmroot/.style={draw=black!75, line width=1.0pt, rounded corners=3.5pt, fill=mmRootFill, inner sep=5pt, font=\normalsize\bfseries},
  mml1/.style={draw=black!55, line width=0.9pt, rounded corners=2.5pt, fill=mmLOneFill, inner sep=4.5pt, font=\small\bfseries, align=left, text width=4.8cm},
  mml2/.style={draw=black!45, line width=0.8pt, rounded corners=2pt, fill=white, inner sep=4pt, font=\small, align=left, text width=6.2cm},
]
  %% 连线：根 → 一级 → 二级（节点填充不透明，盖住线头）
  \draw[black!55, line width=0.9pt] (4.35, 8.3375) -- (4.8, 17.25);
  \draw[black!55, line width=0.9pt] (4.8, 17.25) -- (10.5, 18.975);
  \draw[black!55, line width=0.9pt] (4.8, 17.25) -- (10.5, 17.825);
  \draw[black!55, line width=0.9pt] (4.8, 17.25) -- (10.5, 16.675);
  \draw[black!55, line width=0.9pt] (4.8, 17.25) -- (10.5, 15.525);
  \draw[black!55, line width=0.9pt] (4.35, 8.3375) -- (4.8, 10.925);
  \draw[black!55, line width=0.9pt] (4.8, 10.925) -- (10.5, 14.375);
  \draw[black!55, line width=0.9pt] (4.8, 10.925) -- (10.5, 13.225);
  \draw[black!55, line width=0.9pt] (4.8, 10.925) -- (10.5, 12.075);
  \draw[black!55, line width=0.9pt] (4.8, 10.925) -- (10.5, 10.925);
  \draw[black!55, line width=0.9pt] (4.8, 10.925) -- (10.5, 9.775);
  \draw[black!55, line width=0.9pt] (4.8, 10.925) -- (10.5, 8.625);
  \draw[black!55, line width=0.9pt] (4.8, 10.925) -- (10.5, 7.475);
  \draw[black!55, line width=0.9pt] (4.35, 8.3375) -- (4.8, 4.6);
  \draw[black!55, line width=0.9pt] (4.8, 4.6) -- (10.5, 6.325);
  \draw[black!55, line width=0.9pt] (4.8, 4.6) -- (10.5, 5.175);
  \draw[black!55, line width=0.9pt] (4.8, 4.6) -- (10.5, 4.025);
  \draw[black!55, line width=0.9pt] (4.8, 4.6) -- (10.5, 2.875);
  \draw[black!55, line width=0.9pt] (4.35, 8.3375) -- (4.8, 0.575);
  \draw[black!55, line width=0.9pt] (4.8, 0.575) -- (10.5, 1.725);
  \draw[black!55, line width=0.9pt] (4.8, 0.575) -- (10.5, 0.575);
  \draw[black!55, line width=0.9pt] (4.8, 0.575) -- (10.5, -0.575);
  %% 节点
  \node[mmroot, anchor=east] at (4.35, 8.3375) {2.2 直线及其方程};
  \node[mml1, anchor=west] at (4.6, 17.25) {2.2.1 直线的倾斜角与斜率};
  \node[mml2, anchor=west] at (10.3, 18.975) {直线的倾斜角};
  \node[mml2, anchor=west] at (10.3, 17.825) {斜率的定义};
  \node[mml2, anchor=west] at (10.3, 16.675) {斜率公式};
  \node[mml2, anchor=west] at (10.3, 15.525) {直线的方向向量};
  \node[mml1, anchor=west] at (4.6, 10.925) {2.2.2 直线的方程};
  \node[mml2, anchor=west] at (10.3, 14.375) {直线的点斜式方程};
  \node[mml2, anchor=west] at (10.3, 13.225) {直线在y轴上的截距};
  \node[mml2, anchor=west] at (10.3, 12.075) {直线的斜截式方程};
  \node[mml2, anchor=west] at (10.3, 10.925) {直线的两点式方程};
  \node[mml2, anchor=west] at (10.3, 9.775) {直线的截距式方程};
  \node[mml2, anchor=west] at (10.3, 8.625) {直线的一般式方程};
  \node[mml2, anchor=west] at (10.3, 7.475) {直线系方程};
  \node[mml1, anchor=west] at (4.6, 4.6) {2.2.3 两条直线的位置关系};
  \node[mml2, anchor=west] at (10.3, 6.325) {设两条不重合的直线，，斜率若存在且分别为，，倾斜角分别为，则对应关系如下：};
  \node[mml2, anchor=west] at (10.3, 5.175) {设两条直线，的斜率分别为，，则对应关系如下：};
  \node[mml2, anchor=west] at (10.3, 4.025) {两直线的交点坐标};
  \node[mml2, anchor=west] at (10.3, 2.875) {两直线的位置关系};
  \node[mml1, anchor=west] at (4.6, 0.575) {2.2.4 点到直线的距离};
  \node[mml2, anchor=west] at (10.3, 1.725) {点到直线的距离与两条平行直线间的距离};
  \node[mml2, anchor=west] at (10.3, 0.575) {常见代数式的几何含义};
  \node[mml2, anchor=west] at (10.3, -0.575) {曼哈顿距离};
\end{tikzpicture}
\end{document}
